\chapter{Vergelijking van functies}\label{ch:b2:comparison}

De asymptotische analyse --- de kunst om een ingewikkelde grootheid
te vervangen door een eenvoudige plus een beheerste fout --- begon
in het volume van bachelorjaar 1 met taylorontwikkelingen. Dit
hoofdstuk maakt er een vak op zich van: ontwikkelingen langs
algemene schalen, de vergelijking van reeks en integraal met haar
volle asymptotische kracht, de formule van Stirling (\omterm{def:b2:metric:complete}{volledig}
bewezen), en de systematische studie van impliciet gedefinieerde
rijen. Deze technieken zijn het dagelijks brood van de
asymptotische analyse, en elk later hoofdstuk dat iets schat ---
reeksen, integralen, kansen --- eet mee aan deze tafel.

\section{Vergelijkingsrelaties en schalen}

\begin{definition}\label{def:b2:comparison:landau}
Bij een punt $a$ ($a \in \R$ of $\pm\infty$), voor functies (of
voor rijen, met $n \to \infty$): $f = o(g)$, $f = O(g)$ en $f \sim
g$ als in het volume van bachelorjaar 1. Een
\emph{vergelijkingsschaal}\index{vergelijkingsschaal} bij $a$ is
een familie positieve functies die paarsgewijs vergelijkbaar zijn
en door $o(\cdot)$ totaal geordend worden --- de standaardschaal
bij $+\infty$ is
\[
x^{\alpha} (\ln x)^{\beta}
\qquad (\alpha, \beta \in \R),
\]
lexicografisch geordend naar $(\alpha, \beta)$, en zo nodig
verfijnd met exponentiëlen $\eu^{\gamma x}$.
\end{definition}

\begin{definition}[Asymptotische ontwikkeling]\label{def:b2:comparison:expansion}
$f$ heeft de \emph{asymptotische ontwikkeling}\index{asymptotische
ontwikkeling}
\[
f = c_1 \varphi_1 + c_2\varphi_2 + \dots + c_k \varphi_k +
o(\varphi_k)
\qquad (\varphi_{i+1} = o(\varphi_i) \text{ in de schaal})
\]
wanneer de opeenvolgende resttermen aan de aangegeven schattingen
voldoen. De coëfficiënten liggen dan vast: $c_1 = \lim
f/\varphi_1$, en inductief $c_{i+1} = \lim\,(f - \sum_{j \leq i}
c_j\varphi_j)/\varphi_{i+1}$.
\end{definition}

\begin{example}
Taylorontwikkelingen zijn \omterm{def:b2:comparison:expansion}{asymptotische ontwikkelingen} langs de
schaal $(x - a)^k$ bij $a$. Maar het begrip is strikt ruimer: bij
$+\infty$ geldt
\[
\frac{1}{x - \ln x}
= \frac1x \cdot \frac{1}{1 - \frac{\ln x}{x}}
= \frac1x + \frac{\ln x}{x^2} + o\Bigl(\frac{\ln x}{x^2}\Bigr),
\]
een ontwikkeling langs de gemengde schaal --- geen enkele
stelling van Taylor is hier van toepassing, alleen de meetkundige
ontwikkeling en het rekenen met $o$'s.
\end{example}

\begin{example}[De standaardschaal is werkelijk geordend]\label{ex:b2:comparison:scaleorder}
De lexicografische bewering van \cref{def:b2:comparison:landau}
vraagt per geval één regel bewijs. Vergelijk
$x^{\alpha}(\ln x)^{\beta}$ en $x^{\alpha'}(\ln x)^{\beta'}$ bij
$+\infty$. Is $\alpha < \alpha'$, dan is de verhouding
$x^{\alpha - \alpha'}(\ln x)^{\beta - \beta'} \to 0$, want een
negatieve macht van $x$ verplettert elke macht van $\ln x$ (zet $x
= \eu^t$: $\eu^{(\alpha - \alpha')t}\,t^{\beta - \beta'} \to 0$
volgens de limiet ``exponentieel wint van polynomiaal'' uit het
volume van bachelorjaar 1). Is $\alpha = \alpha'$ en $\beta <
\beta'$, dan is de verhouding rechtstreeks $(\ln x)^{\beta -
\beta'} \to 0$. De lexicografisch geordende paren $(\alpha,
\beta)$ ordenen de schaal dus naar $o(\cdot)$ --- en de
substitutie $x = \eu^t$ is de universele truc voor gemengde
vergelijkingen van machten en logaritmen.
\end{example}

\begin{example}[Een menagerie rangschikken]\label{ex:b2:comparison:ranking}
Schalen moeten \emph{geordend} zijn; hier is de standaardoefening.
Vergelijk bij $+\infty$ de grootheden $n^{10}$,
$\eu^{\sqrt{\ln n}\,\cdot\,\sqrt n}$, $2^n$ en $n^{\ln n}$ door
logaritmen te nemen:
\[
10\ln n
\;\ll\; (\ln n)^2
\;\ll\; \sqrt{n\ln n}
\;\ll\; n\ln 2 ,
\]
waarbij $a_n \ll b_n$ betekent dat $a_n = o(b_n)$; de tweede term
is $\ln(n^{\ln n})$. Exponentiëren behoudt deze strikte kloven
(gaat $\ln u_n - \ln v_n \to -\infty$, dan $u_n/v_n \to 0$), dus
\[
n^{10} = o\bigl(n^{\ln n}\bigr),
\qquad
n^{\ln n} = o\bigl(\eu^{\sqrt{n\ln n}}\bigr),
\qquad
\eu^{\sqrt{n\ln n}} = o(2^n) .
\]
De moraal, tweemaal: vergelijk altijd via logaritmen (verschillen
van logaritmen, niet verhoudingen van logaritmen), en besluit
nooit tot $u_n \sim v_n$ uit $\ln u_n \sim \ln v_n$ --- het paar
$n^{10}$ en $n^{\ln n}$ heeft een $\ln$-verhouding die naar
$\infty$ gaat, terwijl $2^n$ en $4^n$ een $\ln$-verhouding van
precies $2$ hebben en toch wild niet-equivalent zijn.
\end{example}

\section{Vergelijking van reeks en integraal, asymptotisch}

\begin{theorem}\label{thm:b2:comparison:seriesintegral}
Zij $f$ \omterm{def:b2:metric:continuity}{continu}, positief en \emph{dalend} op
$\intco{1}{+\infty}$.
\begin{enumerate}
  \item Convergeert $\int_1^{\infty} f$, dan voldoen de restsommen
        aan
        \[
        \int_{n+1}^{\infty} f \;\leq\; \sum_{k > n} f(k) \;\leq\;
        \int_{n}^{\infty} f .
        \]
  \item Divergeert $\int_1^\infty f$, dan voldoen de partiële
        sommen aan $\sum_{k=1}^{n} f(k) = \int_1^n f + C + o(1)$
        voor zekere constante $C$: het verschil $\sum_{k \leq n}
        f(k) - \int_1^n f$ \emph{convergeert}.
\end{enumerate}
\end{theorem}

\begin{proof}
De insluiting $f(k+1) \leq \int_k^{k+1} f \leq f(k)$ (dalend) was
het middel uit bachelorjaar 1; sommeren over $k \geq n+1$
respectievelijk $k \geq n$ geeft (1). Zet voor (2) $u_k = f(k) -
\int_k^{k+1} f$: uit de insluiting volgt $0 \leq u_k \leq f(k) -
f(k+1)$, dus zijn de partiële sommen van $\sum u_k$ begrensd door
de telescopische som $f(1) - f(n+1) \leq f(1)$: de reeks
convergeert. Bovendien is de rij $\bigl(\int_n^{n+1} f\bigr)_n$
niet-stijgend ($f$ daalt) en niet-negatief, en dus convergent. Met
\[
\sum_{k=1}^{n} f(k) - \int_1^n f
= \sum_{k=1}^{n} u_k + \int_n^{n+1} f
\]
convergeert het rechterlid als $n \to \infty$: het verschil
convergeert naar een constante $C$, en dat is uitspraak (2).
\end{proof}

\begin{example}[De harmonische ontwikkeling]\label{ex:b2:comparison:harmonic}
Voor $f(t) = \frac1t$: $H_n = \ln n + \gamma + o(1)$, waarmee de
constante van Euler (volume van bachelorjaar 1) met een schoner
bewijs wordt teruggevonden. Eén orde verder
(\cref{exo:b2:comparison:3}):
\[
H_n = \ln n + \gamma + \frac{1}{2n} + o\Bigl(\frac1n\Bigr).
\]
Getallen maken de winst zichtbaar bij $n = 10$: $H_{10} =
2.928968\dots$ en $\ln 10 = 2.302585\dots$, dus de ruwe schatting
van $\gamma$ is $H_{10} - \ln 10 = 0.626383$, er $0.049$ naast;
trekken we de correctie $\frac1{20}$ eraf, dan krijgen we
$0.576383$, slechts $8.3\cdot10^{-4}$ van $\gamma = 0.577216$
verwijderd --- en dat is zelf de volgende term
$\frac{1}{12\cdot100}$ van de ontwikkeling, zoals de
weekendopgave bewijst (vraag 8).
\end{example}

\begin{example}[Een ruwe $\ln(n!)$ zonder Stirling]\label{ex:b2:comparison:lnfactorial}
Het insluitingsmiddel alleen lokaliseert $\ln(n!)$ al. Omdat $\ln$
stijgt, geldt
\[
\int_{k-1}^{k}\ln t\,\dd t \;\leq\; \ln k \;\leq\;
\int_{k}^{k+1}\ln t\,\dd t ,
\]
en sommeren over $k = 2, \dots, n$ (met $\int_1^n\ln = n\ln n - n
+ 1$) geeft
\[
n\ln n - n + 1 \;\leq\; \ln(n!) \;\leq\; (n+1)\ln(n+1) - n .
\]
Beide hekken zijn $n\ln n - n + O(\ln n)$: dus $\ln(n!) = n\ln n -
n + O(\ln n)$, en in het bijzonder $\ln(n!) \sim n\ln n$. Wat
Stirling daaraan toevoegt zijn de volgende twee sporten --- de
$\frac12\ln n$ en de constante $\ln\sqrt{2\pi}$ --- die de fijnere
telescopering van \cref{thm:b2:comparison:stirling} kosten. Weten
welke precisie elk gereedschap koopt, is de helft van het
asymptotische vakmanschap.
\end{example}

\begin{example}[Wegvallen vraagt om ontwikkelingen]\label{ex:b2:comparison:cancellation}
Bereken de limiet van $\sqrt{n^2 + n} - n$. Beide termen zijn
$\sim n$, en ``$\sim n - n$'' is betekenisloos: equivalenten
mogen niet worden afgetrokken. Ontwikkel in plaats daarvan:
\[
\sqrt{n^2 + n} - n
= n\Bigl(\sqrt{1 + \tfrac1n} - 1\Bigr)
= n\Bigl(\frac{1}{2n} - \frac{1}{8n^2} +
O\Bigl(\frac1{n^3}\Bigr)\Bigr)
= \frac12 - \frac{1}{8n} + O\Bigl(\frac{1}{n^2}\Bigr) :
\]
de limiet is $\frac12$, met de naderingssnelheid $\frac1{8n}$ als
bonus. Het mechanisme verdient een naam: een verschil van twee
grote equivalente grootheden leeft \omterm{def:b2:metric:complete}{volledig} in hun \emph{volgende}
termen, dus moet men ontwikkelen tot de eerste orde waarin de twee
kanten verschillen --- en de restterm meenemen om te certificeren
dat er in die orde niets anders overleeft.
\end{example}

\begin{example}[Een divergente vergelijking, uitgewerkt]\label{ex:b2:comparison:loglog}
Voor $f(t) = \frac{1}{t\ln t}$ op $\intco{2}{+\infty}$ (\omterm{def:b2:metric:continuity}{continu},
positief, dalend) is $\int_2^x f = \ln\ln x - \ln\ln 2 \to
\infty$, dus geeft \cref{thm:b2:comparison:seriesintegral}~(2)
\[
\sum_{k=2}^{n}\frac{1}{k\ln k} = \ln\ln n + C + o(1)
\]
voor zekere constante $C$. Twee lessen. Ten eerste is de
divergentie echt maar ijzig traag: de partiële som overschrijdt
$4$ pas rond $n \approx \eu^{\eu^{4 - C}}$, astronomisch groot.
Ten tweede werd de \emph{vorm} $\ln\ln n$ geleverd door een
primitieve, niet geraden: voor monotone termen is de integraal het
canonieke somgereedschap, en de constante $C$ --- net als de
$\gamma$ van Euler --- is het geheugen van de eerste termen.
\end{example}

\section{De formule van Stirling}

\begin{lemma}[Wallis-integralen, hernomen]\label{lem:b2:comparison:wallis}
Zij $W_n = \int_0^{\pi/2} \sin^n t\,\dd t$. Dan is $nW_nW_{n-1} =
\frac\pi2$ voor $n \geq 1$, daalt $(W_n)$, en geldt $W_n \sim
\sqrt{\dfrac{\pi}{2n}}$.
\end{lemma}

\begin{proof}
Partiële integratie geeft $nW_n = (n-1)W_{n-2}$ ($n \geq 2$), dus
is $nW_nW_{n-1}$ constant in $n$, gelijk aan $1 \cdot W_1 W_0 =
\frac\pi2$. Dalend: $\sin^{n+1} \leq \sin^n$ op
$\intcc{0}{\frac\pi2}$. De insluiting in detail: de monotonie geeft
$W_{n+1} \leq W_n \leq W_{n-1}$, en deling door $W_{n-1} > 0$
levert
\[
\frac{n}{n+1} = \frac{W_{n+1}}{W_{n-1}} \leq
\frac{W_n}{W_{n-1}} \leq 1 ,
\]
waarbij de linkeridentiteit uit de recurrentie bij index $n + 1$
komt. Beide grenzen gaan naar $1$: dus $W_n \sim W_{n-1}$, waaruit
\[
nW_n^2 \sim nW_nW_{n-1} = \frac\pi2
\qquad\Longrightarrow\qquad
W_n \sim \sqrt{\frac{\pi}{2n}} .
\]
\end{proof}

\begin{example}[De eerste Wallis-integralen]\label{ex:b2:comparison:wallisvalues}
Uit $W_0 = \frac\pi2$, $W_1 = 1$ en de recurrentie $nW_n =
(n-1)W_{n-2}$:
\[
W_2 = \frac\pi4, \qquad
W_3 = \frac23, \qquad
W_4 = \frac{3\pi}{16}, \qquad
W_5 = \frac{8}{15}, \qquad
W_6 = \frac{5\pi}{32}.
\]
De even indices dragen een $\pi$, de oneven zijn rationaal --- de
twee vervlochten producten uit de gesloten vormen. Numeriek is
$W_6 \approx 0.4909$ tegenover de asymptotische waarde
$\sqrt{\pi/12} \approx 0.5116$: bij $n = 6$ zit het equivalent al
binnen $5\%$, en de productidentiteit klopt exact bij elke $n$:
$6\,W_6W_5 = 6\cdot\frac{5\pi}{32}\cdot\frac8{15} = \frac\pi2$.
Kleine tabellen als deze zijn de goedkoopste manier om een
algebrafoutje te vangen voordat het een asymptotisch argument
besmet.
\end{example}

\begin{theorem}[Stirling]\label{thm:b2:comparison:stirling}
\[
n! \;\sim\; \sqrt{2\pi n}\, \Bigl(\frac{n}{\eu}\Bigr)^{\!n}
.\index{formule van Stirling}
\]
\end{theorem}

\begin{proof}
\emph{Stap 1: $n! \sim C \sqrt n\, (n/\eu)^n$ voor zekere
constante $C > 0$.} Zet
\[
d_n = \ln(n!) - \Bigl(n + \frac12\Bigr)\ln n + n .
\]
Dan is
\[
d_n - d_{n+1}
= \Bigl(n + \frac12\Bigr) \ln\frac{n+1}{n} - 1
= \Bigl(n + \frac12\Bigr)\Bigl(\frac1n - \frac{1}{2n^2} +
\frac{1}{3n^3} + o\bigl(n^{-3}\bigr)\Bigr) - 1
= \frac{1}{12n^2} + o\Bigl(\frac{1}{n^2}\Bigr),
\]
volgens de taylorontwikkeling van $\ln(1 + \frac1n)$. De reeks
$\sum (d_n - d_{n+1})$ convergeert dus absoluut (vergelijking met
$\sum n^{-2}$), zodat $(d_n)$ convergeert, zeg naar $d$;
exponentiëren geeft $n! \sim C\sqrt n\,(n/\eu)^n$ met $C =
\eu^{d}$.

\emph{Stap 2: $C = \sqrt{2\pi}$ via Wallis.} De gesloten vorm
$W_{2p} = \frac{(2p)!}{4^p (p!)^2}\cdot\frac\pi2$ (uit de
recurrentie, de berekening uit bachelorjaar 1 opnieuw gedaan in de
context van \cref{lem:b2:comparison:wallis}) combineert met stap
1:
\[
W_{2p} \sim \frac{C\sqrt{2p}\,(2p/\eu)^{2p}}
{4^p\,\bigl(C\sqrt p\,(p/\eu)^p\bigr)^2}\cdot\frac{\pi}{2}
= \frac{\sqrt{2p}}{C\,p}\cdot\frac{\pi}{2}
= \frac{\pi}{C}\cdot\frac{1}{\sqrt{2p}} .
\]
Vergelijking met $W_{2p} \sim \sqrt{\frac{\pi}{4p}}$
(\cref{lem:b2:comparison:wallis}): uit
$\frac{\pi}{C\sqrt{2p}} = \sqrt{\frac{\pi}{4p}}\,(1 + o(1))$ volgt
$C = \pi \sqrt{\frac{4p}{2p\,\pi}} = \sqrt{2\pi}$.
\end{proof}

\begin{example}[Centrale binomiaalcoëfficiënt]\label{ex:b2:comparison:centralbinomial}
\[
\binom{2n}{n} = \frac{(2n)!}{(n!)^2}
\sim \frac{\sqrt{4\pi n}\,(2n/\eu)^{2n}}{2\pi n\,(n/\eu)^{2n}}
= \frac{4^n}{\sqrt{\pi n}} :
\]
de kans dat een symmetrische toevalswandeling op tijdstip $2n$
naar $0$ terugkeert, is $\sim \frac{1}{\sqrt{\pi n}}$ --- een
aankondiging van \cref{ch:b2:randomvar}.
\end{example}

\begin{remark}[Vooruitblik binnen dit volume]\label{rem:b2:comparison:perspectives}
Elk kwantitatief hoofdstuk hierna spreekt de taal van dit
hoofdstuk. \cref{ch:b2:series} klasseert reeksen door de termen
met de schaal $n^{-\alpha}(\ln n)^{-\beta}$ te vergelijken --- de
weekendopgave daar brengt die grens \omterm{def:b2:metric:complete}{volledig} in kaart.
\cref{ch:b2:integration} doet hetzelfde voor oneigenlijke
integralen, met dezelfde schaal in de \omterm{def:b2:metric:continuity}{continue} veranderlijke.
\cref{ch:b2:powerseries} berekent convergentiestralen uit
$\limsup\abs{a_n}^{1/n}$, een oefening in equivalenten van
$n$-de wortels waarbij Stirling de standaardsleutel is
($\sqrt[n]{n!} \sim \frac n\eu$, \cref{exo:b2:comparison:4}). En
de kanshoofdstukken verzilveren Stirling rechtstreeks: de lokale
schattingen van \cref{ch:b2:randomvar} voor
binomiaalcoëfficiënten zijn woordelijk
\cref{ex:b2:comparison:centralbinomial} en
\cref{ex:b2:comparison:lopsided}. De asymptotiek is hier geen
hoofdstuk; zij is het accent van het hele volume.
\end{remark}

\begin{method}[De checklist voor bootstrappen]\label{met:b2:comparison:checklist}
Controleer vier punten voordat je een gebootstrapte ontwikkeling
vertrouwt. (1) \emph{Eerst het bestaan:} het nulpunt of de rij
moet zijn vastgepind (monotonie, tussenwaarden) voordat er ook maar
iets wordt ontwikkeld --- symbolen zonder referent laten zich
prachtig ontwikkelen en betekenen niets. (2) \emph{Eén orde per
ronde:} elke substitutie mag slechts tot de orde van de ingevoerde
schatting worden vertrouwd; twee nieuwe termen uit één ronde
persen is de klassieke bron van verkeerde coëfficiënten. (3)
\emph{Resttermen rijden mee:} sleep de $o(\cdot)$ door elke
algebraïsche stap mee en laat de absorptie (kleinere termen door
grotere resttermen opgeslokt) pas aan het eind gebeuren, en
expliciet. (4) \emph{Numerieke controle:} evalueer bij één
eerlijke waarde van $n$; een fout in een coëfficiënt overleeft een
algebraïsche herafleiding verrassend vaak, en overleeft de
rekenkunde bijna nooit.
\end{method}

\begin{remark}[Klassieke valkuilen]\label{rem:b2:comparison:pitfalls}
(i) Equivalenten \emph{tellen} slecht op: uit $u_n \sim n + \ln n$
en $v_n \sim -n$ mag men \emph{niet} besluiten dat $u_n + v_n \sim
\ln n$; wegvallende termen vragen om ontwikkelingen met expliciete
resttermen, nooit om kale equivalenten. (ii) Exponentieer nooit een
equivalentie: $n + 1 \sim n$ maar $\eu^{n+1} \not\sim \eu^n$; de
veilige richting is logaritmen nemen van equivalenten die naar
$+\infty$ gaan (de weekendopgave van dit hoofdstuk, vraag 24).
(iii) Een \omterm{def:b2:comparison:expansion}{asymptotische ontwikkeling} hoort bij een \emph{schaal}:
schrijven dat $f = \frac1x + o\bigl(\frac1{x^2}\bigr)$ beweert
meer dan $f = \frac1x + o\bigl(\frac1x\bigr)$, en de twee mengen
maakt alle latere algebra ongeldig. (iv) Substitueer bij het
bootstrappen de \emph{hele} huidige ontwikkeling, restterm
inbegrepen --- een $o(\cdot)$ halverwege laten vallen levert
plausibele maar verkeerde coëfficiënten op. (v) De vergelijking
van reeks en integraal heeft monotonie nodig: voor oscillerende
termen faalt zij ronduit (vergelijk $\sum\frac{\sin k}k$,
\cref{ch:b2:series}).
\end{remark}

\begin{example}[Stirling in getallen]\label{ex:b2:comparison:stirlingnum}
Bij $n = 10$ geeft de formule $\sqrt{20\pi}\,(10/\eu)^{10}
\approx 3\,598\,696$ tegenover $10! = 3\,628\,800$: een relatieve
fout van $8.3\cdot10^{-3}$, opmerkelijk voor een
``asymptotische'' uitspraak bij $n = 10$. De fout heeft structuur
--- de exacte verfijning $n! = \sqrt{2\pi n}\,(n/\eu)^n\bigl(1 +
\frac1{12n} + O(n^{-2})\bigr)$ --- en haar eerste correctie
$\frac1{120} \approx 8.3\cdot10^{-3}$ verklaart de waargenomen
kloof vrijwel precies. Het gereedschap van Euler--Maclaurin uit de
weekendopgave is juist de systematische bron van zulke
correctietermen.
\end{example}

\begin{remark}[Waar dit hoofdstuk wordt gebruikt]
De asymptotische vergelijking is de grammatica van alles wat
verderop kwantitatief is: de convergentietests en het
Bertrand-panorama van \cref{ch:b2:series}, de
integreerbaarheidscriteria van \cref{ch:b2:integration}, de
berekeningen van convergentiestralen in
\cref{ch:b2:powerseries}, en de limietstellingen van
\cref{ch:b2:randomvar} (waar Stirling de schattingen van de
Moivre en Laplace aandrijft). Het volume van bachelorjaar 3
industrialiseert het ene idee dat wij hier met de hand bewijzen
--- haal de hoofdterm eruit, begrens de rest --- tot de methode
van Laplace en de gedomineerde convergentie.
\end{remark}

\begin{example}[Een integraal met zichzelf vergeleken: $\int_2^x \frac{\dd t}{\ln t}$]\label{ex:b2:comparison:liworked}
De vergelijkingsgereedschapskist werkt ook op integralen. Zij
$F(x) = \int_2^x\frac{\dd t}{\ln t}$ (de integrand is \omterm{def:b2:metric:continuity}{continu} op
$\intco2\infty$). Partieel integreren:
\[
F(x) = \Bigl[\frac{t}{\ln t}\Bigr]_2^x +
\int_2^x\frac{\dd t}{(\ln t)^2}
= \frac{x}{\ln x} + O\Bigl(\int_2^x\frac{\dd t}{(\ln
t)^2}\Bigr) + O(1),
\]
en de restintegraal is $o\bigl(\frac{x}{\ln x}\bigr)$: splits haar
bij $\sqrt x$ en begrens met
\[
\int_2^{\sqrt x}\frac{\dd t}{(\ln t)^2} \leq \sqrt x
\qquad\text{en}\qquad
\int_{\sqrt x}^{x}\frac{\dd t}{(\ln t)^2} \leq
\frac{x}{(\ln\sqrt x)^2} = \frac{4x}{(\ln x)^2} .
\]
Bijgevolg is $F(x) \sim \frac{x}{\ln x}$. Wie de priemgetalstelling
in de weekendopgave van dit hoofdstuk is tegengekomen, herkent
$F$: het is de logaritmische integraal, de betere schatter van
$\pi(x)$, en de berekening laat zien dat zij tot op de eerste orde
met $\frac{x}{\ln x}$ overeenstemt.
\end{example}

\begin{example}[Stirling op een scheve binomiaalcoëfficiënt]\label{ex:b2:comparison:lopsided}
Dezelfde routine met drie faculteiten als bij
\cref{ex:b2:comparison:centralbinomial} geeft voor
$\binom{3n}{n} = \frac{(3n)!}{n!\,(2n)!}$:
\[
\binom{3n}{n} \sim
\frac{\sqrt{6\pi n}\,(3n/\eu)^{3n}}
{\sqrt{2\pi n}\,(n/\eu)^{n}\cdot\sqrt{4\pi n}\,(2n/\eu)^{2n}}
= \sqrt{\frac{3}{4\pi n}}\,
\Bigl(\frac{27}{4}\Bigr)^{\!n} .
\]
De exponentiële snelheid $\frac{27}4 = \frac{3^3}{2^2}$ is
$\eu^{3n\,H(1/3)}$ in de entropienotatie van de informatietheorie:
scheve binomiaalcoëfficiënten groeien strikt langzamer dan de
centrale $4^n$ per twee stappen --- hier $(27/4)^{1/3} \approx
1.89 < 2$ per stap. Elke asymptotiek van binomiaalcoëfficiënten in
de combinatoriek en de kansrekening (\cref{ch:b2:randomvar}) is
deze ene berekening met andere gewichten.
\end{example}

\section{Impliciet gedefinieerde rijen}

\begin{method}\label{met:b2:comparison:implicit}
Om de asymptotiek van de oplossingen $x_n$ van een vergelijking
$F(x, n) = 0$ te vinden:\index{impliciete rij}
\begin{enumerate}
  \item \emph{Lokaliseer}: bewijs het bestaan en de eenduidigheid
        van $x_n$ in een bepaald interval (monotonie,
        tussenwaardestelling), en bepaal het ruwe gedrag (limiet,
        groeiorde).
  \item \emph{Bootstrap}: substitueer de ruwe vorm $x_n =
        (\text{hoofdterm})(1 + \varepsilon_n)$ in de vergelijking
        en los op naar de volgende orde van $\varepsilon_n$;
        herhaal, waarbij elke ronde één orde verfijnt.
\end{enumerate}
\end{method}

\begin{example}\label{ex:b2:comparison:tan}
Voor $n \geq 1$ heeft de vergelijking $\tan x = x$ precies één
oplossing $x_n$ in $\intoo{n\pi - \frac\pi2}{n\pi + \frac\pi2}$
(de functie $\tan x - x$ stijgt daar van $-\infty$ tot $+\infty$,
want haar afgeleide is $\tan^2 x \geq 0$). \emph{Ruw:} $x_n = n\pi
+ \frac\pi2 - y_n$ met $y_n \in \intoo{0}{\pi}$; omdat $x_n \to
\infty$ en $\tan x_n = x_n \to +\infty$, nadert $x_n$ de asymptoot
van links: $y_n \to 0$. \emph{Bootstrap:} $\tan x_n = \cot y_n =
\frac{1}{\tan y_n} \sim \frac{1}{y_n}$, en de vergelijking $\cot
y_n = x_n \sim n\pi$ geeft $y_n \sim \frac{1}{n\pi}$. Bijgevolg is
\[
x_n = n\pi + \frac\pi2 - \frac{1}{n\pi} +
o\Bigl(\frac1n\Bigr),
\]
en het proces gaat tot elke orde door
(\cref{exo:b2:comparison:6}).
\end{example}

\begin{example}[Een tweede ronde van de methode]\label{ex:b2:comparison:xplusln}
Los $x + \ln x = n$ asymptotisch op. \emph{Lokaliseer:} $x \mapsto
x + \ln x$ stijgt op $\intoo{0}{+\infty}$ van $-\infty$ tot
$+\infty$: dus één nulpunt $x_n$, en $x_n \to \infty$. \emph{Ruw:}
uit $\ln x_n = o(x_n)$ volgt $x_n \sim n$. \emph{Bootstrap:} uit
$x_n = n - \ln x_n$ en $\ln x_n = \ln n + o(1)$ (logaritmen van
equivalenten, beide leden $\to \infty$):
\[
x_n = n - \ln n + o(1) ;
\]
nog een ronde, met $\ln x_n = \ln\bigl(n - \ln n + o(1)\bigr) =
\ln n - \frac{\ln n}{n} + o\bigl(\frac{\ln n}n\bigr)$:
\[
x_n = n - \ln n + \frac{\ln n}{n} +
o\Bigl(\frac{\ln n}{n}\Bigr).
\]
(Controle bij $n = 100$: het nulpunt is $x \approx 95.4415$; de
formule met drie termen geeft $100 - 4.6052 + 0.0461 = 95.4409$ en
die met twee termen $95.3948$ --- elke ronde wint de voorspelde
orde.) Dezelfde lus, een derde landschap: de methode van
\cref{met:b2:comparison:implicit} maalt niet om de gedaante van de
vergelijking, zolang elke ronde de dominante onbekende isoleert.
\end{example}

\section{Oefeningen}

\begin{exercise}[$\star$]\label{exo:b2:comparison:1}
Ontwikkel bij $+\infty$, twee termen voorbij de hoofdterm:
\[
\sqrt{x^2 + x + 1} ,
\qquad
\ln(x^2 + x) - 2\ln x,
\qquad
\frac{x + \sin x}{x - \ln x} .
\]
\end{exercise}

\begin{exercise}[$\star$]\label{exo:b2:comparison:2}
Geef de aard (convergentie of divergentie) en, bij divergentie, de
hoofdasymptotiek van $\sum_{k \leq n} k^\alpha$ voor $\alpha > -1$,
$\alpha = -1$ en $\alpha < -1$, met
\cref{thm:b2:comparison:seriesintegral}.
\end{exercise}

\begin{exercise}[$\star\star$]\label{exo:b2:comparison:3}
Bewijs dat $H_n = \ln n + \gamma + \frac{1}{2n} +
o\bigl(\frac1n\bigr)$. \emph{(Bestudeer $v_n = H_n - \ln n -
\gamma$: toon aan dat $v_n - v_{n+1} = \frac{1}{2n^2} +
O(n^{-3})$ en sommeer de staart, met de vergelijking $\sum_{k \geq
n} \frac{1}{2k^2} \sim \frac{1}{2n}$ ---
\cref{thm:b2:comparison:seriesintegral}~(1).)}
\end{exercise}

\begin{exercise}[$\star\star$]\label{exo:b2:comparison:4}
Bepaal met Stirling equivalenten van $\dfrac{(3n)!}{(n!)^3}$,
$\;\dfrac{n!}{n^n}$ en $\;\sqrt[n]{n!}$ (als $\frac n\eu(1 +
o(1))$, tot op twee termen preciesgemaakt).
\end{exercise}

\begin{exercise}[$\star\star$]\label{exo:b2:comparison:5}
Bewijs voor $n \geq 2$ dat $x^n + x = 1$ precies één oplossing $x_n
\in \intoo{0}{1}$ heeft, dat $x_n \to 1$, en leid af dat
\[
x_n = 1 - \frac{\ln n}{n} + o\Bigl(\frac{\ln n}{n}\Bigr).
\]
\emph{(Uit $x_n^n = 1 - x_n$: neem logaritmen en bootstrap met $x_n
= 1 - \varepsilon_n$.)}
\end{exercise}

\begin{exercise}[$\star\star$]\label{exo:b2:comparison:6}
Voer \cref{ex:b2:comparison:tan} één orde verder door:
\[
x_n = n\pi + \frac\pi2 - \frac{1}{n\pi} + \frac{1}{2n^2\pi} +
o\Bigl(\frac{1}{n^2}\Bigr).
\]
\emph{(Schrijf $\cot y_n = x_n$ exact, ontwikkel $\cot y = \frac1y
- \frac y3 + o(y)$ en $x_n = n\pi(1 + \frac{1}{2n} - \dots)$, en
identificeer.)}
\end{exercise}

\begin{exercise}[$\star\star$]\label{exo:b2:comparison:7}
Bepaal $\lim_{n\to\infty} \dfrac{1}{n!}\sum_{k=0}^{n} k!$
\emph{(begrens de som van alle termen op de laatste twee na)}, en
leid daaruit de \omterm{def:b2:comparison:expansion}{asymptotische ontwikkeling} $\sum_{k \leq n} k! =
n!\bigl(1 + \frac1n + O(n^{-2})\bigr)$ af.
\end{exercise}

\begin{exercise}[$\star\star\star$]\label{exo:b2:comparison:8}
Zij $u_0 > 0$ en $u_{n+1} = u_n + \dfrac{1}{u_n}$. Bewijs dat $u_n
\to \infty$, vervolgens dat $u_n \sim \sqrt{2n}$ \emph{(bestudeer
$u_n^2$: haar toenamen zijn $2 + u_n^{-2}$; sommeer)}, en verfijn:
\[
u_n = \sqrt{2n}\Bigl(1 + \frac{\ln n}{8n} +
o\Bigl(\frac{\ln n}{n}\Bigr)\Bigr).
\]
\emph{(Uit $u_n^2 = 2n + \sum_{k<n} u_k^{-2} + u_0^2$ en $u_k^2
\sim 2k$: de som is $\sim \frac12\ln n$ volgens
\cref{thm:b2:comparison:seriesintegral}.)}
\end{exercise}

\begin{exercise}[$\star\star\star$]\label{exo:b2:comparison:9}
(Een riemannsom met een addertje) Bepaal het asymptotische gedrag
van
\[
S_n = \sum_{k=1}^{n} \frac{1}{n + k\ln n} .
\]
\emph{(Haal $n$ buiten haakjes: $S_n = \frac1n\sum_k \bigl(1 +
\frac{k\ln n}{n}\bigr)^{-1}$; herken een som van het type Riemann
met een traag variërende parameter $t = \ln n$, bereken $\int_0^1
\frac{\dd u}{1 + tu} = \frac{\ln(1+t)}{t}$, en besluit dat $S_n
\sim \frac{\ln\ln n}{\ln n}$.)}
\end{exercise}

\begin{exercise}[$\star$]\label{exo:b2:comparison:10}
Bewijs de identiteit $(\ln n)^{\ln n} = n^{\ln\ln n}$, en
rangschik daarna de volgende grootheden bij oneindig in stijgende
$o(\cdot)$-volgorde, met bewijs: $n^2$, $(\ln n)^{\ln n}$, $2^n$,
$n!$, $n^n$.
\end{exercise}

\begin{exercise}[$\star\star$]\label{exo:b2:comparison:11}
(Staart van $\sum 1/k^2$, twee termen) Bewijs met de exacte
telescopering $\sum_{k > n} \frac{1}{k(k+1)} = \frac{1}{n+1}$ en de
splitsing $\frac1{k^2} = \frac{1}{k(k+1)} + \frac{1}{k^2(k+1)}$
dat
\[
\sum_{k > n} \frac{1}{k^2}
= \frac1n - \frac{1}{2n^2} + O\Bigl(\frac{1}{n^3}\Bigr).
\]
\end{exercise}

\begin{exercise}[$\star\star\star$]\label{exo:b2:comparison:12}
Zij $u_0 = \frac12$ en $u_{n+1} = u_n + \eu^{-u_n}$. Bewijs dat
$u_n \to \infty$, en leid vervolgens --- door $v_n = \eu^{u_n}$ te
zetten en aan te tonen dat $v_{n+1} = v_n + 1 + \frac{1}{2v_n} +
O\bigl(v_n^{-2}\bigr)$ --- af dat
\[
u_n = \ln n + \frac{\ln n}{2n} + O\Bigl(\frac1n\Bigr).
\]
\end{exercise}

\section{Probleem: Bootstrappen, van Euler--Maclaurin tot de
priemgetallen}

Een impliciete of opgestapelde grootheid geeft haar asymptotiek
zelden in één keer prijs; men haalt haar in rondes binnen, waarbij
elke ronde de vorige schatting terugvoert in de definiërende
betrekking. Deze weekendopgave traint die lus op nieuwe
vergelijkingen, bewijst de \emph{formule van Euler--Maclaurin} van
eerste orde (de trapeziumverbetering van de vergelijking van reeks
en integraal, met rigoureuze foutbalken), keert $x\ln x = n$ om, en
int de beroemdste cheque van de methode: uit de zonder bewijs
aangenomen priemgetalstelling volgt de asymptotische wet $p_n \sim
n\ln n$ voor het $n$-de priemgetal.

\begin{problem}[{Weekendopgave --- de correctie van
Euler--Maclaurin en de asymptotiek van het $n$-de priemgetal}]
\label{pb:b2:comparison:1}

\medskip
\noindent\textbf{Deel I --- De bootstraplus op een nieuwe
vergelijking.}
\begin{enumerate}
  \item Bewijs de bewering over eenduidigheid uit
        \cref{def:b2:comparison:expansion}: geldt $f = \sum_{i\leq
        k} c_i\varphi_i + o(\varphi_k) = \sum_{i \leq k}
        c_i'\varphi_i + o(\varphi_k)$ langs dezelfde schaal, dan
        is $c_i = c_i'$ voor alle $i$. Voer daarna het gemengde
        voorbeeld uit de cursus één sport verder door:
        \[
        \frac{1}{x - \ln x} = \frac1x + \frac{\ln x}{x^2} +
        \frac{(\ln x)^2}{x^3} + o\Bigl(\frac{(\ln
        x)^2}{x^3}\Bigr) \qquad (x \to +\infty),
        \]
        en leg uit waarom er geen term $\frac{c}{x^2}$
        verschijnt.
  \item Toon aan dat de vergelijking $\eu^x + x = n$ voor elke $n
        \geq 1$ precies één reële oplossing $x_n$ heeft, en dat
        $x_n \to +\infty$ met $x_n \sim \ln n$.
  \item Bootstrap tweemaal:
        \[
        x_n = \ln n - \frac{\ln n}{n} - \frac{(\ln n)^2}{2n^2} +
        o\Bigl(\frac{(\ln n)^2}{n^2}\Bigr).
        \]
  \item Controleer numeriek bij $n = 1000$: vergelijk $x_{1000}
        \approx 6.90083$ met de waarden van vraag 3 met één, twee
        en drie termen, tot op vijf decimalen.
\end{enumerate}

\medskip
\noindent\textbf{Deel II --- Euler--Maclaurin, eerste orde.}
\begin{enumerate}[resume]
  \item Bewijs de identiteit met de trapeziumkern: voor $g$ van
        klasse $C^2$ op $\intcc{0}{1}$ geldt
        \[
        \int_0^1 g(t)\,\dd t = \frac{g(0) + g(1)}{2}
        - \frac12\int_0^1 t(1 - t)\,g''(t)\,\dd t
        \]
        \emph{(integreer $\frac12 t(1-t)g''$ tweemaal partieel)}.
  \item Zij $f$ van klasse $C^2$ op $\intco{1}{+\infty}$ met
        $\int_1^\infty \abs{f''} < \infty$. Toon aan dat
        \[
        E_n = \sum_{k=1}^{n} f(k) - \int_1^n f -
        \frac{f(1) + f(n)}{2}
        \]
        naar een constante $E$ convergeert, met de staartgrens
        $\abs{E - E_n} \leq \frac18\int_n^\infty\abs{f''}$: de
        \emph{formule van Euler--Maclaurin} tot op de eerste
        orde.
  \item Pas dit toe op $f(t) = \frac1t$: bewijs dat
        \[
        H_n = \ln n + \gamma + \frac{1}{2n} + \varepsilon_n,
        \qquad \abs{\varepsilon_n} \leq \frac{1}{8n^2},
        \]
        wat \cref{exo:b2:comparison:3} versterkt (identificeer de
        constante met $\gamma$ door met
        \cref{ex:b2:comparison:harmonic} te vergelijken).
  \item Haal de volgende coëfficiënt binnen: toon aan dat
        $\varepsilon_n = -\frac{1}{12n^2} +
        o\bigl(\frac1{n^2}\bigr)$ \emph{(de toenamen van $E_n$
        zijn $\frac12\int_0^1t(1-t)f''(n+t)\dd t =
        \frac1{12}f''(n) + o(f''(n))$; sommeer de staart met
        \cref{thm:b2:comparison:seriesintegral})}.
  \item Pas vraag 6 toe op $f = \ln$: leid in drie regels opnieuw
        de convergentie af van $d_n = \ln n! - (n + \frac12)\ln n
        + n$ (stap 1 van \cref{thm:b2:comparison:stirling}), met
        als bonus de foutsnelheid $d_n = d +
        O\bigl(\frac1n\bigr)$.
  \item Pas vraag 6 toe op $f(t) = \frac{1}{\sqrt t}$: toon aan
        dat
        \[
        \sum_{k=1}^{n}\frac1{\sqrt k} = 2\sqrt n + c +
        \frac{1}{2\sqrt n} + O\Bigl(\frac{1}{n^{3/2}}\Bigr)
        \]
        voor zekere constante $c$, en evalueer alle termen bij $n
        = 10^4$ (de constante is $c \approx -1.4604$).
\end{enumerate}

\medskip
\noindent\textbf{Deel III --- Omkeren: de vergelijking $x\ln x =
n$.}
\begin{enumerate}[resume]
  \item Toon aan dat $x\ln x = n$ voor $n \geq 1$ precies één
        oplossing $x_n \in \intco{1}{+\infty}$ heeft, dat $x_n \to
        \infty$, en dat $\ln x_n \sim \ln n$.
  \item Leid de omkering met één term af, $x_n \sim
        \dfrac{n}{\ln n}$, en bootstrap daarna nog eens:
        \[
        \ln x_n = \ln n - \ln\ln n + o(1),
        \qquad
        x_n = \frac{n}{\ln n}\Bigl(1 + \frac{\ln\ln n}{\ln n}
        + o\Bigl(\frac{\ln\ln n}{\ln n}\Bigr)\Bigr).
        \]
  \item Test bij $n = 10^6$: het werkelijke nulpunt is $x \approx
        87\,848$; vergelijk met de waarden met één term ($\approx
        72\,382$) en met twee termen ($\approx 86\,140$), en
        verklaar de trage winst (de ontwikkelingsparameter is
        $\frac{\ln\ln n}{\ln n}$, bij $n = 10^6$ slechts $\approx
        0.19$).
  \item We \emph{nemen} nu de priemgetalstelling zonder bewijs
        \emph{aan}: het aantal $\pi(x)$ priemgetallen $\leq x$
        voldoet aan $\pi(x) \sim \frac{x}{\ln x}$ als $x \to
        \infty$ (eerlijk bewezen in het volume van bachelorjaar
        3). Schrijf $p_n$ voor het $n$-de priemgetal, verantwoord
        dat $\pi(p_n) = n$, en voer de omkering van de vragen
        11--12 uit om te bewijzen dat
        \[
        p_n \sim n \ln n .
        \]
  \item Opbrengsten: (a) toon aan dat $\sum_{k \leq n} p_k \sim
        \frac{n^2\ln n}{2}$ \emph{(vergelijk $\sum k\ln k$ met
        $\int t\ln t\,\dd t$)}; (b) bereken bij benadering de kans
        dat een uniform gekozen geheel getal van $100$ cijfers
        priem is ($\ln 10^{100} \approx 230.26$: ongeveer één op
        $230$).
\end{enumerate}

\medskip
\noindent\textbf{Deel IV --- De methode geëxporteerd: $x\tan x =
1$.}
\begin{enumerate}[resume]
  \item Toon aan dat de vergelijking $\tan x = \frac1x$ voor elke
        $n \geq 1$ precies één oplossing $x_n$ heeft in
        $\intoo{n\pi}{\,n\pi + \frac\pi2}$, en dat $z_n = x_n -
        n\pi \to 0^+$.
  \item Eén term: $z_n \sim \dfrac{1}{n\pi}$.
  \item Toon aan dat de ontwikkeling van $z_n$ \emph{geen} term
        $\frac{c}{n^2}$ heeft: $z_n = \frac1{n\pi} +
        O\bigl(\frac{1}{n^3}\bigr)$.
  \item Drie termen: bewijs met $\arctan u = u - \frac{u^3}3 +
        O(u^5)$ en $\frac1{x_n} = \frac{1}{n\pi} -
        \frac{z_n}{(n\pi)^2} + O(n^{-3}\cdot z_n^2)$ dat
        \[
        x_n = n\pi + \frac{1}{n\pi} -
        \frac{4}{3\pi^3 n^3} + o\Bigl(\frac{1}{n^3}\Bigr).
        \]
  \item Controleer bij $n = 3$: het werkelijke nulpunt is $x_3
        \approx 9.5293344$; vergelijk de waarden met één en met
        drie termen, en zet dit in één zin af tegen de $\tan x =
        x$ uit de cursus (\cref{ex:b2:comparison:tan}): waar elke
        rij in haar venster ligt, en waarom.
\end{enumerate}

\medskip
\noindent\textbf{Deel V --- Een dynamische bootstrap, de
spelregels, en de synthese.}
\begin{enumerate}[resume]
  \item Zij $u_0 \in \intoo{0}{\pi}$ en $u_{n+1} = \sin u_n$. Toon
        aan dat $u_n$ dalend naar $0$ gaat, en bereken de limiet
        van $\dfrac{1}{u_{n+1}^2} - \dfrac{1}{u_n^2}$
        \emph{(ontwikkel $\sin^{-2}$ via $\sin u = u -
        \frac{u^3}6 + o(u^3)$)}.
  \item Leid daaruit, via de gemiddelden van Cesàro (volume van
        bachelorjaar 1), de klassieke uitspraak
        \[
        u_n \sim \sqrt{\frac{3}{n}}
        \]
        af.
  \item (Gecertificeerde numeriek) Toon met de rigoureuze grens
        van vraag 7 aan dat het evalueren van $\ln n + \gamma +
        \frac1{2n}$ bij $n = 10^6$ de waarde $H_{10^6}$ oplevert
        met een fout van hoogstens $1.25\cdot10^{-13}$ --- een som
        van een miljoen termen, met drie termen tot op dertien
        cijfers berekend.
  \item (Spelregels) Bewijs of weerleg, met bewijzen of
        tegenvoorbeelden: (a) geldt $u_n \sim v_n \to +\infty$,
        dan $\ln u_n \sim \ln v_n$; (b) geldt $u_n \sim v_n$, dan
        $\eu^{u_n} \sim \eu^{v_n}$; (c) geldt $f \sim g$ bij
        $+\infty$ (met $f, g$ differentieerbaar), dan $f' \sim
        g'$.
  \item (Synthese) In telkens één zin: de bootstraplus van
        \cref{met:b2:comparison:implicit} zoals gebruikt in de
        Delen I, III en IV; wat de trapeziumcorrectie toevoegt aan
        \cref{thm:b2:comparison:seriesintegral}; waarom het
        omkeren van $x\ln x$ juist de brug is van $\pi(x)$ naar
        $p_n$; en welke van de regels uit vraag 24 welke stap
        beschermde. Noem de twee toppen: de formule van
        Euler--Maclaurin (eerste orde), en de asymptotische wet
        voor het $n$-de priemgetal.
\end{enumerate}
\end{problem}
